\chapter{KẾ HOẠCH XÂY DỰNG, KIỂM THỬ VÀ ĐÁNH GIÁ}

\section{Cách tổ chức công việc}

Ba bài kiểm tra thường xuyên và bài thi cuối được thực hiện trên cùng một sản phẩm, không tách thành các báo cáo độc lập. Các yêu cầu phân tích--thiết kế, phần quản lý, phần tích hợp AI và phần hoàn thiện cuối kỳ được dùng như những cổng kiểm tra nối tiếp. Nhóm sẽ xây dựng theo từng phần chạy được; sau mỗi phần sẽ kiểm tra và ghi lại kết quả trước khi mở rộng. Cách này phù hợp với tình trạng hiện tại của dự án: tài liệu đã có nhưng mã nguồn chưa bắt đầu.

Trong cùng một báo cáo, bốn mốc được hiểu theo phạm vi sau: bài kiểm tra thường xuyên 1 tập trung vào phân tích bài toán và thiết kế; bài kiểm tra thường xuyên 2 tập trung vào cấu trúc dự án và phần quản lý cốt lõi; bài kiểm tra thường xuyên 3 tập trung vào tích hợp, kiểm thử và kiểm soát chức năng AI; bài thi cuối kiểm tra bản chạy hoàn chỉnh, chất lượng dữ liệu và giao diện, bảo mật, hiệu năng, triển khai, báo cáo và thuyết trình. Các mốc này chỉ là cách kiểm tra tiến độ của một sản phẩm, không phải bốn phiên bản tài liệu tách rời.

Một chức năng chỉ được xem là sẵn sàng khi có mô tả nghiệp vụ, thiết kế liên quan, mã nguồn và kiểm thử. Hình chụp giao diện chỉ dùng để minh họa cách sử dụng, không thay thế cho kiểm thử hoặc mã nguồn.

\section{Cấu trúc mã nguồn dự kiến}

\begin{table}[h!]
\centering
\small
\caption{Cấu trúc thư mục dự kiến}
\label{tab:v8-repository}
\begin{tabularx}{\textwidth}{L{4.0cm}Y Y}
\toprule
\textbf{Thư mục} & \textbf{Nội dung} & \textbf{Tài liệu/kết quả đi kèm}\\
\midrule
\texttt{backend/app/api} & Đường dẫn API và dữ liệu vào/ra. & Tài liệu API và kiểm thử API.\\
\texttt{backend/app/domain} & Quy tắc chiến dịch, nội dung, chỉ số và quyền. & Kiểm thử nghiệp vụ và lịch sử quyết định.\\
\texttt{backend/app/ai} & Chọn dữ liệu, prompt, gọi provider và kiểm tra kết quả. & Prompt, nhật ký gọi AI và kiểm thử lỗi.\\
\texttt{backend/tests} & Kiểm thử đơn vị, tích hợp và API. & Báo cáo chạy test và dữ liệu kiểm thử.\\
\texttt{frontend/src} & Màn hình, biểu mẫu, tìm kiếm, dashboard và trạng thái lỗi. & Kịch bản demo và kiểm tra giao diện.\\
\texttt{db/migrations} và \texttt{db/seed} & Thay đổi schema và dữ liệu mẫu. & Cách tạo lại cơ sở dữ liệu.\\
\texttt{docs} & Hướng dẫn chạy, bảo mật, API và đánh giá. & README, runbook và báo cáo.\\
\bottomrule
\end{tabularx}
\end{table}

\section{Các giai đoạn thực hiện}

\begin{table}[h!]
\centering
\small
\caption{Các giai đoạn từ thiết kế đến bản trình diễn}
\label{tab:v8-roadmap}
\begin{tabularx}{\textwidth}{L{1.5cm}L{3.2cm}Y Y}
\toprule
\textbf{Giai đoạn} & \textbf{Mục tiêu} & \textbf{Phạm vi} & \textbf{Điều kiện chuyển bước}\\
\midrule
1 & Nền tảng & Đăng nhập, vai trò, sản phẩm, kênh và chiến dịch. & CRUD, kiểm tra quyền và dữ liệu mẫu chạy ổn định.\\
2 & Nội dung & Bản nháp, tìm kiếm, lịch đăng, trạng thái và phê duyệt. & Có kiểm tra chuyển trạng thái đúng/sai.\\
3 & Chỉ số & Nhập metrics, tính KPI và dashboard. & Công thức khớp dữ liệu mẫu, xử lý đủ dữ liệu rỗng.\\
4 & AI & Gợi ý, tạo bản nháp, tóm tắt và nhật ký. & Có kiểm tra cấu trúc, quyền, cảnh báo và lỗi provider.\\
5 & Hoàn thiện & Bảo mật, kiểm thử lại, hướng dẫn chạy, backup và demo. & Người khác có thể làm theo hướng dẫn hoặc blocker được ghi rõ.\\
\bottomrule
\end{tabularx}
\end{table}

\section{Kế hoạch kiểm thử}

\begin{table}[h!]
\centering
\scriptsize
\caption{Các nhóm kiểm thử bắt buộc}
\label{tab:v8-test-plan}
\begin{tabularx}{\textwidth}{L{1.0cm}L{4.0cm}Y Y}
\toprule
\textbf{Mã} & \textbf{Trường hợp} & \textbf{Kết quả cần có} & \textbf{Dạng kiểm thử}\\
\midrule
T01 & Đăng nhập đúng/sai và truy cập theo vai trò. & Phiên đúng được tạo; phiên sai và thao tác trái quyền bị từ chối. & API và quyền.\\
T02 & Tạo chiến dịch thiếu trường, ngân sách âm hoặc ngày sai. & Không tạo bản ghi; lỗi chỉ rõ trường sai. & API và nghiệp vụ.\\
T03 & Khóa ngoại không tồn tại hoặc bản ghi trùng. & Giao dịch bị hủy; không có dữ liệu mồ côi. & Cơ sở dữ liệu.\\
T04 & Tìm kiếm, lọc và sắp xếp chiến dịch. & Kết quả đúng với dữ liệu mẫu và bộ lọc. & Tích hợp.\\
T05 & Gửi duyệt, từ chối, chấp thuận và lập lịch. & Chỉ chuyển trạng thái hợp lệ; từ chối có lý do. & Quy trình và giao diện.\\
T06 & Nhập metrics, trùng ngày/kênh và mẫu số bằng 0. & KPI đúng; không chia cho 0 hoặc nhân đôi số liệu. & Dữ liệu và truy vấn.\\
T07 & AI nhận dữ liệu ngoài quyền hoặc context quá dài. & Dữ liệu bị loại hoặc rút gọn có thông báo. & Quyền và chọn dữ liệu.\\
T08 & AI trả JSON đúng/sai cấu trúc. & Chỉ JSON hợp lệ mới tạo bản nháp dùng được. & Hợp đồng dữ liệu.\\
T09 & AI timeout, giới hạn tốc độ hoặc provider ngừng hoạt động. & Có thời hạn, thử lại giới hạn và giữ bản nháp cũ. & Giả lập lỗi.\\
T10 & Kết quả AI có thông tin không được cung cấp. & Có cảnh báo; không tự chấp thuận hoặc xuất bản. & Kiểm tra nội dung.\\
T11 & Log chứa token, mật khẩu hoặc thông tin thừa. & Bí mật bị loại hoặc che; không trả ra giao diện. & Bảo mật.\\
T12 & Cài đặt mới, seed, backup và khôi phục. & Làm lại được theo README; sai khác được ghi lại. & Tái lập.\\
T13 & Người dùng hoàn thành luồng tạo chiến dịch và gửi nội dung duyệt. & Giao diện có loading/empty/error state; thao tác sai được giải thích. & Kịch bản sử dụng.\\
T14 & Đo truy vấn dashboard và gọi AI trên fixture cố định. & Ghi p50/p95, error rate, môi trường và kích thước dữ liệu; chưa có ngưỡng trước khi đo. & Đo hiệu năng.\\
T15 & Người khác dựng lại dự án từ repository sạch. & README, migration/seed và backup/restore cho kết quả có thể đối chiếu. & Tái lập.\\
\bottomrule
\end{tabularx}
\end{table}

\section{Cách đánh giá}

Nhóm sẽ đánh giá ba khía cạnh riêng và đối chiếu chúng với các tiêu chí của bài thi cuối:

\begin{enumerate}
\item \textbf{Tính đúng của phần quản lý:} tỷ lệ kiểm thử đạt, tính toàn vẹn của dữ liệu và việc thực hiện đúng quyền/trạng thái.
\item \textbf{Chất lượng chức năng AI:} tỷ lệ kết quả đúng cấu trúc, mức bám sát dữ liệu, mức phù hợp với mục tiêu/kênh, cảnh báo và lượng chỉnh sửa của người dùng.
\item \textbf{Khả năng sử dụng lại:} thời gian cài đặt, dữ liệu mẫu, phiên bản mã nguồn, môi trường, câu lệnh và kết quả chạy lại.
\end{enumerate}

Với AI, một bộ trường hợp cố định sẽ được dùng cho ít nhất ba lần thử. Mỗi lần chỉ thay đổi prompt, model hoặc tham số đang so sánh. Kết quả tốt và kết quả lỗi đều phải được lưu; không chỉ chọn những câu trả lời đẹp để đưa vào báo cáo. Nếu chưa kết nối được provider thật, nhóm chỉ đánh giá được quy trình và phần kiểm tra kết quả, không kết luận về chất lượng mô hình.

\section{Phương án triển khai}

Hình~\ref{fig:v8-deployment} là phương án triển khai nhỏ nhất cho prototype. Đây là thiết kế để chuẩn bị môi trường, chưa phải mô tả của hệ thống đã chạy.

\begin{figure}[h!]
\centering
\begin{tikzpicture}[
  font=\sffamily\footnotesize,
  >=Latex,
  node_box/.style={
    draw=ictunavy,
    fill=white,
    rounded corners=4pt,
    minimum width=3.6cm,
    minimum height=1.7cm,
    align=center,
    line width=0.85pt,
    text width=3.4cm,
    inner sep=4pt
  },
  cloud_box/.style={
    draw=black!70,
    fill=black!2,
    rounded corners=4pt,
    minimum width=3.3cm,
    minimum height=1.6cm,
    align=center,
    line width=0.8pt,
    dashed,
    text width=3.1cm,
    inner sep=4pt
  },
  rel/.style={draw=black!85,line width=0.75pt,-{Latex[length=2.2mm]}},
  comm/.style={font=\sffamily\tiny,fill=white,inner sep=1.5pt,text=black,align=center}
]

% Cột 1: Client Node (x = 1.9cm, y = 0)
\node[node_box] (client) at (1.9, 0) {
  \textbf{$\ll$device$\gg$ Máy trạm Client}\\[2pt]
  {\tiny [Web Browser / PC / Mobile]}\\[4pt]
  \fbox{\tiny React SPA Bundle}
};

% Cột 2: Máy chủ Backend (x = 7.0cm, y = 1.5) & Cơ sở dữ liệu (x = 7.0cm, y = -1.5)
\node[node_box] (server) at (7.0, 1.5) {
  \textbf{$\ll$device$\gg$ Máy chủ Backend}\\[2pt]
  {\tiny [Docker Container / Linux]}\\[4pt]
  \fbox{\tiny FastAPI Engine / Python}
};

\node[node_box] (database) at (7.0, -1.5) {
  \textbf{$\ll$device$\gg$ Máy chủ CSDL}\\[2pt]
  {\tiny [Dedicated Storage Server]}\\[4pt]
  \fbox{\tiny PostgreSQL / SQLite}
};

% Cột 3: Dịch vụ ngoài (x = 12.1cm)
\node[cloud_box] (ai_cloud) at (12.1, 1.5) {
  \textbf{$\ll$cloud service$\gg$ Dịch vụ AI}\\[2pt]
  {\tiny [OpenAI / Gemini API]}\\[2pt]
  {\tiny Xử lý prompt, sinh JSON nháp}
};

\node[cloud_box] (soc_cloud) at (12.1, -1.5) {
  \textbf{$\ll$cloud service$\gg$ Mạng Xã hội}\\[2pt]
  {\tiny [Meta / TikTok / Google API]}\\[2pt]
  {\tiny Xuất bản bài viết sau duyệt}
};

% Kết nối rõ ràng, thẳng hàng, không chéo đè chữ
% 1. Client -> Server
\draw[rel] (client.east) -- node[comm,above,sloped]{HTTPS / TLS 1.3} (server.west);

% 2. Server -> Database (kết nối dọc)
\draw[rel] (server.south) -- node[comm,right=2pt]{TCP 5432 / SQL} (database.north);

% 3. Server -> AI Cloud (kết nối ngang)
\draw[rel] (server.east) -- node[comm,above=2pt]{HTTPS / JSON} (ai_cloud.west);

% 4. Server -> Social Cloud (kết nối xiên)
\draw[rel] (server.south east) -- node[comm,above,sloped]{OAuth 2.0 / REST} (soc_cloud.north west);

\end{tikzpicture}

\caption{Phương án triển khai prototype}
\label{fig:v8-deployment}
\end{figure}

README cuối kỳ cần ghi phiên bản Python/Node, cách tạo môi trường, biến trong \texttt{.env.example}, cách chạy migration/seed, lệnh khởi động frontend/backend, tài khoản demo, lệnh chạy kiểm thử và cách backup/restore. Key thật không được đưa vào mã nguồn, ảnh chụp hoặc nhật ký.

\section{Kịch bản trình diễn}

\begin{enumerate}
\item Quản lý đăng nhập, tạo sản phẩm và tạo một chiến dịch mùa hè với mục tiêu, đối tượng, ngân sách và thời gian.
\item Quản lý thêm kênh Facebook và Email, phân công nhân viên; nhân viên mở chiến dịch và yêu cầu gợi ý ý tưởng.
\item Nhân viên chọn một ý tưởng, tạo bản nháp theo kênh, xem nguồn dữ liệu và sửa nội dung.
\item Nhân viên gửi duyệt. Quản lý xem, chấp thuận hoặc từ chối có lý do; chỉ nội dung được chấp thuận mới được lập lịch.
\item Người dùng nhập số liệu mẫu theo ngày/kênh và xem CTR, CVR, CPC trên dashboard.
\item Quản lý yêu cầu tóm tắt. Kịch bản cũng phải cho thấy trường hợp provider lỗi hoặc kết quả sai cấu trúc mà không làm mất dữ liệu cũ.
\end{enumerate}

\section{Rủi ro cần theo dõi}

\begin{table}[h!]
\centering
\scriptsize
\caption{Rủi ro và hướng xử lý}
\label{tab:v8-risks}
\begin{tabularx}{\textwidth}{L{3.4cm}L{2.0cm}Y}
\toprule
\textbf{Rủi ro} & \textbf{Mức} & \textbf{Hướng xử lý}\\
\midrule
Phần AI mở rộng quá nhiều & Cao & Giữ ba chức năng trong đề bài; hoàn thiện phần quản lý trước.\\
Provider không ổn định hoặc hết quota & Cao & Có lớp adapter, dữ liệu giả lập và xử lý lỗi; không để AI chặn phần CRUD.\\
Thiếu dữ liệu thật & Trung bình & Dùng dữ liệu mẫu có ghi rõ; chỉ kết luận về quy trình và công thức.\\
SQLite không phù hợp khi nhiều người dùng & Trung bình & Giữ lớp truy cập dữ liệu độc lập; đánh giá PostgreSQL trước khi mở rộng.\\
Thiếu thời gian kiểm thử & Cao & Kiểm thử ngay sau từng giai đoạn, ưu tiên quyền, trạng thái, dữ liệu và lỗi AI.\\
\bottomrule
\end{tabularx}
\end{table}

\section{Kết luận chương}

Kế hoạch này biến tiêu chí của học phần thành những phần việc và phép kiểm tra cụ thể. Sau khi có mã nguồn, các ô kế hoạch phải được thay bằng kết quả thật, kèm môi trường và cách chạy; không giữ nguyên các con số dự kiến như thể đã đo.
